\section{Discussion and Conclusion} \label{sec:8discussion-conclusion}

\noindent\textbf{Threats to validity.}
Our methodology rests on three explicit assumptions that readers should bear in mind when interpreting the verification results.
First, we assume that all keys are uniformly distributed over their value space, so that the Shannon entropy and min-entropy of a key coincide with $\log_2 N_K$.
Protocols in which keys are drawn from non-uniform distributions---for example, user-chosen passwords whose effective entropy is much lower than their nominal length---may be over-approximated as secure by our current model; extending \texttt{MDerive} with non-uniform entropy annotations is an explicit direction for future work.
Second, the oracle-activation budget $B$ is a termination heuristic rather than a soundness condition: a verification that succeeds at budget $B$ rules out attacks using at most $B$ oracle invocations per target term, but a larger budget could in principle surface additional attacks, so reported bounds should be read as ``secure under the stated budget.''
Third, our case studies, while substantial, cover three protocol families and six BC device configurations; generalizing empirical claims about scalability and automation to arbitrary protocols requires further evaluation.

\noindent\textbf{Verification effort.}
Although we impose a budget on the value-breaking oracle, verification may still fail to terminate or take an excessively long time on densely interleaved multi-session models; the worst-case BC configuration in our evaluation takes roughly $9.9 \times 10^4$ seconds.
Compared with baseline models that do not carry low-entropy annotations, the analysis cost grows substantially because the relationship-breaking oracle introduces many new derivation branches.
In practice, difficult BC configurations still benefit from \Tamarin{}'s oracle-heuristic and from a small number of auxiliary lemmas that prune search space; we disclose each case where such human guidance is applied in Section~\ref{sec:6results}.
Engineering further automation---for example, learned heuristics for oracle scheduling or automatic auxiliary-lemma synthesis---is an orthogonal line of work that we see as directly benefiting our approach.

\noindent\textbf{Implications for the Bluetooth Classic specification.}
The UMD attack is not an implementation bug but a weakness in the protocol design: as long as two endpoints can repeatedly renegotiate the session-key length asymmetrically, and as long as the resulting session keys remain derivable fragments of a shared long-term secret, the attack complexity $C(n) = n \cdot 2^{128/n}$ applies regardless of which device vendor is in use.
Based on the attack, we argue that the Bluetooth SIG should consider the following revisions.
(i) Impose a \emph{hard floor} on the session-key length in the core specification (e.g., 16~bytes for Secure, and at least 7~bytes for Enhanced devices communicating with modern peers), so that the $n = 16$ optimum of $C(n)$ becomes unreachable.
(ii) Require \emph{bilateral confirmation} of any key-length downgrade, so that an attacker coercing one side into a short-key session cannot proceed without simultaneously coercing the other.
(iii) Bind short-key negotiations to the pair of \texttt{BDADDR}s that initiated the long-term key, so that fragments obtained in one session cannot be recombined into a key that targets an unrelated session.
We have shared these recommendations with the Bluetooth SIG as part of responsible disclosure.

\noindent\textbf{Scope of the methodology.}
Our methodology currently does not model private (secret) functions, because the prerequisites for brute-force attacks are defined in terms of publicly-known cryptographic operators.
When a real protocol uses a private function, one can often replace it with a combination of public functions and private variables, and this replacement is sound with respect to our attacker model; we leave native support for private functions to future work.
Similarly, our modeling focuses on brute-force attacks whose cost depends solely on key length and does not cover side-channel-assisted or differential attacks whose cost depends on implementation-level details.

\noindent\textbf{Conclusion.}
We proposed the first protocol-agnostic symbolic methodology for analyzing protocols that tolerate low-entropy keys, built on the \texttt{MDerive} derivation abstraction and a three-tier breaking-oracle design, and realized as EntropyVerif---a syntax-sugar extension of \Tamarin{} whose verification soundness is inherited from the unmodified \Tamarin{} backend.
Our case studies on WPA2-PSK, BLE pairing, and Bluetooth Classic automatically rediscover PMKID, KNOB, and Method Confusion without any protocol-specific reveal rules, and reveal a new protocol-level attack on Bluetooth Classic---the Unilateral Multi-Downgrade (UMD) attack (Section~\ref{sec:umd-attack})---that, via the compound-key decomposition $C(n) = n \cdot 2^{128/n}$ attained at $n=16$, reduces 128-bit session-key recovery by a factor of $2^{116}$.
Together, these results demonstrate that entropy-aware symbolic verification is both feasible and necessary for modern protocols that compromise on key strength for usability or compatibility, and they argue for treating brute-force reachability as a first-class citizen in future protocol design and analysis.
